% ID: SEL-D0012-BB-E03
% ESTADO: cotejada. Segunda pasada visual completa y borrador compilado revisados.
% Fuente: Selectividad/matematicas_aplicadas/EXTRA_RESERVA_B_Modelo-6-Enunciados (Mat-Ap-2021).pdf
% SHA-256: 1752ec2f46b4ef2cc5d1c3346d69b79c2cc8a7aa12f6a0cd8f0d994b5bcda157
% Página: 1; bloque B; ejercicio original 3.

Se considera la función
\[
f(x)=\begin{cases}
ax+b & \text{si }x<1\\
x^2-bx+a & \text{si }x\geq 1
\end{cases}
\]
\begin{enumerate}[label=\alph*)]
  \item \textbf{(0.5 puntos)} Halle el valor de $b$ para que $f$ sea continua
  en $\mathbb{R}$.
  \item \textbf{(0.5 puntos)} Para $b=\dfrac{1}{2}$, halle el valor de $a$
  para que $f$ sea derivable en $\mathbb{R}$.
  \item \textbf{(0.7 puntos)} Para $a<0$ y $b=\dfrac{1}{2}$, estudie el
  crecimiento y halle las abscisas de los extremos de la función $f$.
  \item \textbf{(0.8 puntos)} Para $a=0$ y $b=\dfrac{1}{2}$, represente
  la región del plano delimitada por la gráfica de $f$, el eje de abscisas
  y las rectas $x=0$ y $x=2$. Calcule el área de dicha región.
\end{enumerate}
